\part*{Prefacio}

Tomé estos apuntes cuando cursé Diseño y Análisis de Algoritmos en la Universidad de los Andes, entre agosto y noviembre de 2022. 

Posteriormente los fui corrigiendo y extendiendo intermitentemente a medida que trabajaba en temas relacionados con algoritmia. Las extensiones más notables corresponden a dos momentos: cuando cursé Teoría de Grafos entre agosto y noviembre de 2023 y cuando fui seleccionado para dictar el curso de Diseño de Algoritmos en mayo de 2026.

Los apuntes están presentados en el mismo formato en el que originalmente fueron escritos: la clase de \LaTeX{} que establece su estilo la escribí en noviembre del 2020. Más importante aún, los apuntes preservan la estructura y el tono de cuando los escribí como estudiante, con la intención de que sean útiles precisamente a esa audiencia.

La selección deliberada de la audiencia es en gran medida lo que otorga valor a estas notas. Textos clásicos como CLRS\footnote{Uso esta sigla para referirme al famosísimo Introduction to Algorithms de Cormen, Leiserson, Rivest y Stein, que es un excelente ejemplo de este fenómeno: con docenas de miles de citas en google scholar, claramente tiene la profundidad suficiente para resultar útil a investigadores.} intentan dirigirse a una audiencia mixta: estudiantes, profesores, investigadores y profesionales. Esa amplitud de audiencia genera que el texto se difunda ampliamente, pero también implica un nivel de formalidad, profundidad y exhaustividad que muchas veces no es el más adecuado para un primer encuentro con el tema, especialmente a nivel de pregrado.

Estas notas no sacrifican rigor ni correctitud, pero sí priorizan claridad pedagógica sobre exhaustividad formal. En consecuencia, omiten algunos detalles cuya inclusión sería valiosa en un texto de referencia, pero que para muchos estudiantes solo alargarían el texto y añadirían complejidad innecesaria.
% Dónde se evidencia eso?
% En merge sort, prefiero slicing que pasar idxs por todo lado
% En max subarreglo con divide conquer, CLRS pone if-else menos legibles pq quieren algo más parecido a una prueba matemática; en SW de producción, uno querría un early return